\chapter{Le paludisme et la fièvre Ebola}

\noindent\textit{Le paludisme et la fièvre Ebola sont deux maladies infectieuses graves,
particulièrement présentes en Afrique centrale et de l'Ouest. Leurs agents responsables,
leurs modes de transmission et leurs moyens de lutte sont très différents : les
connaître permet de mieux s'en protéger.}
\vspace{8pt}

\connecteBanner

\section{Le paludisme}

\begin{definition}[Paludisme]
Le \textbf{paludisme} est une maladie causée par un parasite unicellulaire du genre
\textbf{\textit{Plasmodium}}, transmis à l'être humain par la piqûre d'un
\textbf{moustique} appelé \textbf{anophèle} femelle.
\end{definition}

\begin{propriete}[Le cycle de la maladie]
En piquant une personne infectée, l'anophèle femelle absorbe le parasite avec le sang.
En piquant ensuite une autre personne, elle lui injecte le parasite. Celui-ci se
multiplie d'abord dans le foie, puis envahit et détruit les \textbf{globules
rouges}.
\end{propriete}

\begin{illus}
\begin{tikzpicture}[scale=0.85]
\draw[onbgreen,fill=onbgreen,fill opacity=0.12,line width=1.2pt] (0,0) circle (0.85);
\node[font=\scriptsize,align=center] at (0,0) {moustique\\infecté};
\draw[->,onbmuted,line width=1.1pt] (0.95,0.3) to[out=30,in=150] (2.9,0.3);
\node[above,font=\scriptsize,onbmuted] at (1.9,0.65) {piqûre};
\draw[onbo,fill=onbo,fill opacity=0.12,line width=1.2pt] (3.8,0) circle (0.85);
\node[font=\scriptsize,align=center] at (3.8,0) {personne\\infectée};
\draw[->,onbmuted,line width=1.1pt] (2.9,-0.3) to[out=-150,in=-30] (0.95,-0.3);
\node[below,font=\scriptsize,onbmuted] at (1.9,-0.65) {nouvelle piqûre};
\end{tikzpicture}
\fig{Figure — Cycle simplifié de transmission du paludisme entre moustique et
personnes.}
\end{illus}

\begin{propriete}[Symptômes]
Fièvre élevée, frissons, maux de tête, courbatures, fatigue intense. Dans les formes
graves (souvent chez le jeune enfant ou la femme enceinte), le paludisme peut
provoquer une \textbf{anémie} sévère, une atteinte du cerveau (paludisme cérébral),
voire le décès en l'absence de traitement rapide.
\end{propriete}

\begin{propriete}[Moyens de lutte]
\begin{itemize}
\item dormir sous une \textbf{moustiquaire imprégnée} d'insecticide ;
\item éliminer les \textbf{eaux stagnantes} autour des habitations (lieux de
reproduction des moustiques) ;
\item utiliser des \textbf{insecticides} et des répulsifs ;
\item consulter rapidement en cas de fièvre pour un \textbf{diagnostic} et un
\textbf{traitement antipaludéen} précoces.
\end{itemize}
\end{propriete}

\begin{remarque}
Les personnes porteuses saines de la drépanocytose (génotype AS, voir chapitre 3)
présentent une résistance partiellement accrue au paludisme : c'est l'une des raisons
pour lesquelles l'allèle S reste fréquent dans certaines régions où le paludisme est
très présent.
\end{remarque}

\section{La fièvre Ebola}

\begin{definition}[Fièvre Ebola]
La \textbf{fièvre Ebola} est une maladie grave causée par un \textbf{virus},
transmise par \textbf{contact direct} avec le sang, les sécrétions ou d'autres
fluides corporels d'une personne (ou d'un animal) infecté. Elle ne se transmet
\textbf{pas} par l'air.
\end{definition}

\begin{propriete}[Symptômes]
Fièvre soudaine et intense, fatigue extrême, douleurs musculaires, maux de tête, mal
de gorge, suivis de vomissements, de diarrhées, et, dans les formes les plus graves,
d'\textbf{hémorragies} internes ou externes.
\end{propriete}

\begin{propriete}[Moyens de lutte]
\begin{itemize}
\item \textbf{isolement rapide} des personnes malades, pour éviter tout nouveau
contact ;
\item port d'équipements de protection individuelle (gants, masques, combinaisons)
par le personnel soignant ;
\item \textbf{désinfection} rigoureuse des lieux et objets contaminés ;
\item identification et suivi des personnes ayant été en contact avec un malade ;
\item la \textbf{vaccination}, disponible depuis quelques années, qui protège les
personnes les plus exposées.
\end{itemize}
\end{propriete}

\begin{remarque}
Piège fréquent : contrairement au paludisme, la fièvre Ebola ne se transmet pas par
un vecteur comme le moustique, mais par un \textbf{contact direct} avec les fluides
corporels d'une personne infectée. Les gestes de lutte contre les deux maladies sont
donc très différents.
\end{remarque}

\begin{propriete}[Comparer paludisme et fièvre Ebola]
\begin{center}
\small
\begin{tabular}{lcc}
\toprule
& \textbf{Paludisme} & \textbf{Fièvre Ebola} \\
\midrule
Agent & parasite (\textit{Plasmodium}) & virus \\
Transmission & piqûre de moustique & contact direct (fluides) \\
Prévention clé & moustiquaire imprégnée & isolement, protection \\
\bottomrule
\end{tabular}
\end{center}
\end{propriete}

% ═══════════════════════════════════════════════════════════════
\section{L'essentiel à retenir}

\begin{synthese}
\textbf{Paludisme.} Parasite \textit{Plasmodium}, transmis par la piqûre de
l'anophèle femelle ; détruit les globules rouges.\\[4pt]
\textbf{Lutte contre le paludisme.} Moustiquaire imprégnée, élimination des eaux
stagnantes, traitement précoce.\\[4pt]
\textbf{Fièvre Ebola.} Virus transmis par contact direct avec les fluides corporels
d'une personne infectée, non par l'air.\\[4pt]
\textbf{Lutte contre Ebola.} Isolement des malades, protection du personnel
soignant, désinfection, vaccination.\\[4pt]
\textbf{Piège fréquent.} Le paludisme se transmet par un vecteur (moustique) ;
Ebola se transmet par contact direct, pas par un vecteur ni par l'air.
\end{synthese}

% ═══════════════════════════════════════════════════════════════
\section{Méthodes \& savoir-faire}

\methode{Distinguer paludisme et fièvre Ebola}
Identifier le mode de transmission (piqûre de moustique ou contact direct avec des
fluides corporels) pour reconnaître la maladie et le geste de prévention adapté.
\begin{exemple}
Une fièvre apparue après plusieurs piqûres de moustique dans une zone à risque évoque
un paludisme plutôt qu'une fièvre Ebola.
\end{exemple}

% ═══════════════════════════════════════════════════════════════
\section{Exercices d'application}
\begin{multicols}{2}

\rubrique{Le paludisme}
\exo{}\diff{1} Quel est l'agent responsable du paludisme ?

\exo{}\diff{1} Par quel insecte le paludisme est-il transmis ?

\exo{}\diff{2} Quelles cellules sanguines sont détruites lors du paludisme ?

\exo{}\diff{1} Citer un moyen de lutte contre le paludisme.

\rubrique{La fièvre Ebola}
\exo{}\diff{1} Quel type d'agent cause la fièvre Ebola ?

\exo{}\diff{2} Comment se transmet la fièvre Ebola ?

\exo{}\diff{1} Citer un moyen de lutte contre la fièvre Ebola.

\rubrique{Comparer les deux maladies}
\exo{}\diff{2} Quelle est la principale différence entre les modes de transmission du
paludisme et de la fièvre Ebola ?

% ═══════════════════════════════════════════════════════════════
\end{multicols}

\section{Exercices d'entraînement}
\begin{multicols}{2}

\rubrique{Le paludisme}
\exo{}\diff{3} Expliquer pourquoi éliminer les eaux stagnantes autour des habitations
contribue à lutter contre le paludisme.

\exo{}\diff{2} Pourquoi les jeunes enfants sont-ils particulièrement vulnérables au
paludisme grave ?

\exo{}\diff{3} Expliquer pourquoi les personnes porteuses saines de la drépanocytose
(AS) sont partiellement protégées contre le paludisme.

\rubrique{La fièvre Ebola}
\exo{}\diff{3} Expliquer pourquoi l'isolement rapide des malades est une mesure
essentielle de lutte contre la fièvre Ebola.

\exo{}\diff{2} Pourquoi le personnel soignant porte-t-il des équipements de
protection individuelle face à un patient atteint d'Ebola ?

\rubrique{Comparer les deux maladies}
\exo{}\diff{3} Un élève affirme qu'une moustiquaire protège aussi contre la fièvre
Ebola. Corriger cette affirmation.

% ═══════════════════════════════════════════════════════════════
\end{multicols}

\section{Sujets type BEPC}

\sujetbac[Le paludisme]
\begin{enumerate}[label=\arabic*)]
\item Quel est l'agent responsable du paludisme et comment est-il transmis ?
\item Cite deux symptômes du paludisme.
\item Cite deux moyens de lutte contre le paludisme.
\item Explique pourquoi la lutte contre les eaux stagnantes limite la transmission
du paludisme.
\end{enumerate}

\sujetbac[La fièvre Ebola]
\begin{enumerate}[label=\arabic*)]
\item Quel type d'agent cause la fièvre Ebola ?
\item Comment cette maladie se transmet-elle ?
\item Cite deux mesures de lutte contre la fièvre Ebola.
\item Explique la différence essentielle entre la transmission du paludisme et celle
de la fièvre Ebola.
\end{enumerate}

% ═══════════════════════════════════════════════════════════════
\section{Corrigés}
\begin{multicols}{2}

\corrige{de l'exercice 1}
Un parasite unicellulaire du genre \textit{Plasmodium}.

\corrige{de l'exercice 2}
Par la piqûre d'un moustique anophèle femelle.

\corrige{de l'exercice 3}
Les globules rouges.

\corrige{de l'exercice 4}
Par exemple : dormir sous une moustiquaire imprégnée.

\corrige{de l'exercice 5}
Un virus.

\corrige{de l'exercice 6}
Par contact direct avec le sang, les sécrétions ou d'autres fluides corporels d'une
personne infectée.

\corrige{de l'exercice 7}
Par exemple : l'isolement rapide des malades.

\corrige{de l'exercice 8}
Le paludisme se transmet par un vecteur (le moustique) ; la fièvre Ebola se transmet
par contact direct avec les fluides corporels, sans vecteur.

\corrige{de l'exercice 9}
Car les eaux stagnantes constituent le lieu de reproduction des moustiques : les
supprimer réduit la population de moustiques et donc le risque de piqûres
infectantes.

\corrige{de l'exercice 10}
Car leur système immunitaire est encore immature et ils ont eu moins d'occasions de
développer une résistance partielle au parasite, ce qui les expose davantage aux
formes graves.

\corrige{de l'exercice 11}
Car chez les personnes AS, la présence de l'allèle S modifie légèrement les globules
rouges, ce qui gêne le développement du parasite \textit{Plasmodium} à l'intérieur de
ces cellules.

\corrige{de l'exercice 12}
Car cela évite tout nouveau contact entre le malade et d'autres personnes, coupant
ainsi la chaîne de transmission du virus.

\corrige{de l'exercice 13}
Pour éviter tout contact direct avec le sang ou les fluides corporels du patient,
seule voie de transmission du virus.

\corrige{de l'exercice 14}
C'est faux : la moustiquaire protège contre les piqûres de moustique (paludisme),
mais la fièvre Ebola ne se transmet pas par les moustiques ; elle se transmet par
contact direct avec des fluides corporels.

\corrige{du sujet 1}
1) Un parasite du genre \textit{Plasmodium}, transmis par la piqûre d'un moustique
anophèle femelle.\\
2) Par exemple : fièvre élevée et frissons.\\
3) Par exemple : moustiquaire imprégnée et élimination des eaux stagnantes.\\
4) Car les eaux stagnantes sont le lieu de reproduction des moustiques ; les
supprimer réduit leur nombre et donc les piqûres infectantes.

\corrige{du sujet 2}
1) Un virus.\\
2) Par contact direct avec le sang, les sécrétions ou d'autres fluides corporels
d'une personne infectée.\\
3) Par exemple : l'isolement des malades et le port d'équipements de protection par
le personnel soignant.\\
4) Le paludisme se transmet par un vecteur (le moustique) ; la fièvre Ebola se
transmet par contact direct, sans vecteur ni transmission par l'air.

\end{multicols}
\exercicesBanner
